\paragraph{Preguntas para investigar}

\begin{enumerate}[label=P\arabic*]
	\item\label{item:mcd}
		?`Qu\'e es el \emph{m\'aximo com\'un divisor} (MCD) de dos enteros?
	\item\label{item:farey}
		?`Qu\'e es una \emph{sucesión de Farey} de orden $3$?
		?`De orden $7$? ?`Y de orden $n$? ¿Cu\'antos elementos hay en una sucesi\'on de
		Farey de orden $3$? ?`Y en una de orden $7$? ?`Y de orden $11$?
		?`Pod\'es hacer alguna conjetura? Representar la fracci\'on $\frac m n$ con el punto
		$(m,n)$ en el plano. Dibujar todas las entradas de la sucesi\'on de Farey de orden $7$.
	\seti
\end{enumerate}
%

\paragraph{Una pregunta para explorar}
\begin{enumerate}[label=P\arabic*]
	\conti
	\item\label{item:cuadratica}
		Encontrar todas las soluciones a la ecuaci\'on $x^2-6x+8=0$ en $\Enterosmod[15]$.
	\seti
\end{enumerate}
%

Probar o refutar las siguientes afirmaciones, y arreglar, si es posible.
\begin{enumerate}[label=P\arabic*]
	\conti
	\item\label{item:divisibilidad:combinacion}
		Si $d\mid a$ y $d\mid b$, y $x,y\in\Enteros$, entonces $d\mid (ax+by)$.
	\item\label{item:divisibilidad:orden}
		Si $a\mid b$ entonces $a\leq b$.
	\item\label{item:divisibilidad:negativa}
		Si $a\nmid b$ y $a\nmid c$, entonces $a\nmid bc$.
	\item\label{item:divisibilidad:ocho}
		Si $2\nmid n$ entonces $8\mid (n^2-1)$.
	\seti
\end{enumerate}
%

\paragraph{Problemas num\'ericos}
\begin{enumerate}[label=P\arabic*]
	\conti
	\item\label{item:algoritmo:euclides}
		Usar el algoritmo de Euclides para hallar el MCD de $29$ y $11$.
	\item\label{item:algoritmo:continua:racional}
		Usar el procedimiento de \ref{item:algoritmo:euclides}
		para probar que:
    		\begin{displaymath}
			\frac{29}{11}\,=\,2+\frac 1 {1+\frac 1 {1+\frac 1 {1+\frac 1 {3}}}}
			\text{ .}
		\end{displaymath}
		%
    		Esta es la \emph{fracci\'on continua simple} de $\frac{29}{11}$.
	\item\label{item:algoritmo:continua:real}
		Cualquier n\'umero real puede ser aproximado por una fracci\'on
		$\frac{a}{b}$ con cualquier grado de precisi\'on, si permitimos que el denominador
		$b$ sea suficientemente grande.
		El tama\~no del denominador es el precio que debemos pagar por una buena aproximaci\'on
		?`Cu\'al es el precio que pagamos aproximando $\frac{29}{11}$ por las siguientes fracciones?
    		\begin{displaymath}
			2\quad\text{,}\quad
			2+\frac{1}{1}\quad\text{,}\quad
			2+\frac{1}{1+\frac{1}{1}}\quad\text{,}\quad
			2+\frac{1}{1+\frac{1}{1+\frac{1}{1}}}
			\text{ .}
		\end{displaymath}
		%
	\item\label{item:unidades}
		Escribir todos los elementos de $\Unidadesmod[7]$, $\Unidadesmod[15]$, $\Unidadesmod[18]$ y
		$\Unidadesmod[21]$ ?`Cu\'ales de ellos son c\'{\i}clicos? ?`Ten\'es alguna conjetura?
	\item\label{item:logaritmos}
		Construir una tabla de ``logaritmos'' para $\Unidadesmod[17]$.
	\item\label{item:logaritmos:ecuaciones}
		Usar la tabla de \ref{item:logaritmos} para resolver las siguientes
		ecuaciones en $\Enterosmod[17]$:
		\begin{displaymath}
			x^2\,=\,2\quad\text{,}\quad
			7x^2\,=\,6\quad\text{,}\quad
			x^3\,=\,3\text{ .}
		\end{displaymath}
		%
		% \begin{enumerate}[label=(\alph*)]
			% \item\label{item:logaritmos:cuadraticos}
				% $x^2=2$;
			% \item\label{item:logaritmos:cuadraticos:bis}
				% $7x^2=6$;
			% \item\label{item:logaritmos:cubicos}
				% $x^3=3$.
		% \end{enumerate}
		% %
	\item\label{item:mcd:ejemplos}
		Encontrar el MCD de cada uno de los siguientes pares de enteros:
		\begin{displaymath}
			(6,15)\quad\text{,}\quad
			(36,49)\quad\text{,}\quad
			(483,291)\quad\text{,}\quad
			(11413,11289)\text{ .}
		\end{displaymath}
		%
		% \begin{enumerate}[label=(\alph*)]
			% \item\label{item:mcd:ejemplos:i}
				% $6,\;15$;
			% \item\label{item:mcd:ejemplos:ii}
				% $36,\;49$;
			% \item\label{item:mcd:ejemplos:iii}
				% $483,\;291$;
			% \item\label{item:mcd:ejemplos:iv}
				% $11413,\;11289$.
		% \end{enumerate}
		% %
		Comparar la eficiencia de cada uno de los siguientes tres m\'etodos a medida que los enteros
		involucrados crecen.
		\begin{enumerate}[label=(\arabic*)]
			\item\label{item:mcd:metodos:divisores}
				Escribir todos los divisores de cada entero y elegir, de entre los divisores
				comunes, el m\'aximo.
			\item\label{item:mcd:metodos:factores}
				Factorizar cada entero como producto de n\'umeros primos y luego hallar el MCD.
			\item\label{item:mcd:metodos:euclides}
				Usar el algoritmo de Euclides.
		\end{enumerate}
		%
	\seti
\end{enumerate}
%

\paragraph{Ingenio}
\begin{enumerate}[label=P\arabic*]
	\conti
	\item\label{item:bases}
		Hallar tres d\'{\i}gitos no nulos $a_0$, $a_1$ y $a_2$, y una base $b$ tal que
		$(a_0a_1a_2)_b=2\,(a_0a_1a_2)_{10}$.
	\item\label{item:bases:fracciones}
		La suma $1+\frac{1}{3}+\frac{1}{5}+\cdots+\frac{1}{2n-1}$ para $n>1$ nunca es un entero.
	% \seti
% \end{enumerate}
% %
% 
% % Agrego uno sencillito pero que puede estar divertido (de tipo ``explorativo''):
% \begin{enumerate}[label=P\arabic*]
	% \conti
	% \item\label{item:extras}
		% Se marcan en una circunferencia $8$ puntos distintos.
		% Probar que se pueden formar la misma cantidad de tri\'angulos que de pent\'agonos
		% usando dichos puntos cómo v\'ertices.
\end{enumerate}
%

